% Generated draft: requires mathematical and source review.
\documentclass[11pt]{article}
\usepackage[margin=25mm]{geometry}
\usepackage{fontspec}
\setmainfont{Latin Modern Roman}
\usepackage{amsmath,amssymb,mathtools,unicode-math}
\setmathfont{Latin Modern Math}
\usepackage{enumitem,needspace,xurl,hyperref}
\setlength{\parindent}{0pt}
\setlength{\parskip}{5pt}
\setlength{\emergencystretch}{3em}
\newcommand{\N}{\mathbb{N}}
\newcommand{\Z}{\mathbb{Z}}
\newcommand{\Q}{\mathbb{Q}}
\newcommand{\R}{\mathbb{R}}
\newcommand{\C}{\mathbb{C}}
\newcommand{\F}{\mathbb{F}}
\newcommand{\A}{\mathbb{A}}
\newcommand{\PP}{\mathbb{P}}
\newcommand{\E}{\mathbb{E}}
\newcommand{\eps}{\varepsilon}
\newcommand{\dd}{\,\mathrm{d}}
\begin{document}
\section*{UnsolvedMath ID: 801113}
\textbf{Source problem code:} catalogue entry
\subsection{Definitions and mathematical statement}
A \(\mathrm{II}_1\) factor is a von Neumann algebra \(N\) with center \(\C 1\), no nonzero minimal projections, and a faithful normal tracial state \(\tau_N\) satisfying \(\tau_N(1)=1\). Use the usual language of tracial von Neumann algebras, with variables ranging over specified operator-norm balls and embeddings understood to be unital, normal, injective, and trace preserving.

Under the standard meaning of existential closedness, \(N\) is existentially closed if, for every embedding \(\iota:N\to M\) into a \(\mathrm{II}_1\) factor, every finite parameter tuple \(\bar a\) from \(N\), and every real-valued quantifier-free formula \(\varphi(\bar x,\bar y)\), one has
\[
\inf_{\bar x\in N}\varphi^N(\bar x,\bar a)
=
\inf_{\bar x\in M}\varphi^M(\bar x,\iota(\bar a)),
\]
where each coordinate of \(\bar x\) ranges over the ball specified by its sort. Thus every such embedding is existential. The supplied wording says that whenever \(N\) embeds into \(M\), it ``can embed existentially'' into \(M\); it does not explicitly distinguish the standard condition on every given embedding from the weaker-looking assertion that some existential embedding exists. The questions below use the standard meaning, subject to this qualification. No separability hypothesis on the factors is supplied.

For a discrete group \(G\), let \(\lambda_G\) denote its left regular representation on \(\ell^2(G)\), defined by
\[
\lambda_G(g)\delta_h=\delta_{gh}\qquad(g,h\in G),
\]
and let
\[
L(G)=\{\lambda_G(g):g\in G\}'',\qquad
\tau_G(x)=\langle x\delta_e,\delta_e\rangle.
\]
Here the double commutant is taken in \(B(\ell^2(G))\), and \(e\) is the identity of \(G\).

\begin{enumerate}
\item Does there exist a countable discrete group \(G\) such that \(L(G)\), with its canonical trace, is an existentially closed \(\mathrm{II}_1\) factor?
\item Does there exist an existentially closed \(\mathrm{II}_1\) factor \(N\) which is not isomorphic, as a von Neumann algebra, to a group von Neumann algebra?
\end{enumerate}
In the second question, the source does not specify whether the comparison groups must be countable. The two corresponding precise targets are
\[
\exists N\ \bigl[N\text{ is an existentially closed }\mathrm{II}_1\text{ factor and }
\forall H\text{ discrete},\ N\not\cong L(H)\bigr]
\]
and the variant in which \(H\) ranges only over countable discrete groups. This draft preserves that unresolved scope distinction rather than identifying the two formulations. The source attributes the existence of existentially closed factors to reference [ECfactors] and points to [IOANA2024110264] and [GKEPTConj] for concrete constructions; complete references are not supplied.
\subsection{Short English statement}
The first question asks whether the group von Neumann algebra of some countable group can be an existentially closed \(\mathrm{II}_1\) factor. The second asks whether an existentially closed \(\mathrm{II}_1\) factor can fail to be a group von Neumann algebra. The source leaves the countability scope of the second question and the precise embedding quantifier in its introductory definition unstated.
\subsection{Sources}
\begin{itemize}
\item Local UnsolvedMath registry, permanent ID \(801113\), source problem code \texttt{catalogue entry}, title \emph{Ultraproduct \(\mathrm{II}_1\) Factors --- Problem 14}. The supplied record attributes the problem to Elayavalli, proposed in 2025, and labels it open. This status is supplied registry data, not an independently verified mathematical assessment.
\item Elayavalli, \emph{50 Open Problems: Ultraproduct \(\mathrm{II}_1\) Factors} (2025), Problem 14. Supplied source URL: \url{https://arxiv.org/abs/2511.20377}. The record reports access on 2026-10-16 and extraction by \texttt{arxiv\_latex}; these provenance details are supplied and were not independently verified here.
\item Reference [ECfactors], cited in the supplied statement for existence of existentially closed factors. No bibliographic details or separate URL are supplied for this reference.
\item References [IOANA2024110264] and [GKEPTConj], cited in the supplied statement for concrete constructions. No bibliographic details or separate URLs are supplied for these references.
\item Supplied upstream revision: \texttt{d371ee917ac9647fbb06cd6fc529deec6f3dda82}. The registry record reports creation and update on 2026-10-16.
\end{itemize}
\end{document}
